\chapter{الأعداد العشرية}\label{ch:g6:decimals}

بين $2$ و $3$ أعداد كثيرة: $2.5$ و $2.71$
و $2.999$\,\dots\ \omterm{def:g6:decimals:places}{الأعداد العشرية} تمدّ نظام المراتب إلى
يمين الآحاد، وتتبع قواعد الأعداد الصحيحة نفسها ---
مع فخاخ قليلة سيعلّمك هذا الفصل كيف تتجنبها.

\section{الكتابة العشرية}

\begin{definition}[المراتب العشرية]\label{def:g6:decimals:places}
\emph{العدد العشري}\index{عدد عشري} له جزء صحيح وجزء
عشري، تفصل بينهما \omterm{def:g4:decimals:point}{النقطة العشرية}. وكل مرتبة على يمين
النقطة تساوي عُشر ما قبلها: أعشار، وأجزاء من مئة،
وأجزاء من ألف. فمثلًا
\[
13.407 = 13 + \frac{4}{10} + \frac{0}{100} + \frac{7}{1000}
= 1 \times 10 + 3 + 4 \times \frac{1}{10} + 7 \times \frac{1}{1000}.
\]
\end{definition}

\begin{omfigure}
\begin{tikzpicture}[scale=0.9]
  \draw (0,0) grid[xstep=2.3, ystep=0.85] (11.5,1.7);
  \node[font=\small] at (1.15,1.275) {عشرات};
  \node[font=\small] at (3.45,1.275) {آحاد};
  \node[font=\small] at (5.75,1.275) {أعشار};
  \node[font=\small] at (8.05,1.275) {أجزاء من مئة};
  \node[font=\small] at (10.35,1.275) {أجزاء من ألف};
  \node[font=\large, omDef] at (1.15,0.425) {$1$};
  \node[font=\large, omDef] at (3.45,0.425) {$3$};
  \node[font=\large, omThm] at (5.75,0.425) {$4$};
  \node[font=\large, omThm] at (8.05,0.425) {$0$};
  \node[font=\large, omThm] at (10.35,0.425) {$7$};
  \fill (4.6,0.06) circle (2.5pt);
\end{tikzpicture}

{\small جدول المراتب للعدد $13.407$: \omterm{def:g4:decimals:point}{النقطة العشرية} تقع
بين الآحاد والأعشار.}
\end{omfigure}

\begin{remark}[أصفار مهمّة وأصفار غير مهمّة]
إضافة الأصفار في \emph{آخر} الجزء العشري لا تغيّر شيئًا:
$2.5 = 2.50 = 2.500$. أما الأصفار \emph{بين} الأرقام فضرورية:
$13.407 \neq 13.47$. والعدد $0.5$ مختلف جدًّا عن $0.05$!
\end{remark}

\begin{method}[مقارنة الأعداد العشرية]\label{met:g6:decimals:compare}
\begin{enumerate}
  \item قارن الأجزاء الصحيحة أولًا: $7.2 > 5.99$ لأن $7 > 5$.
  \item وإن تساوت الأجزاء الصحيحة، فقارن الأجزاء العشرية رقمًا
        رقمًا من اليسار --- والحشو بأصفار نهائية يساعد:
        لمقارنة $3.4$ و $3.15$، اُكتب $3.40$ و $3.15$؛
        وبما أن $40 > 15$ جزءًا من مئة، فإن $3.4 > 3.15$.
\end{enumerate}
واِنتبه: \emph{الأطول لا يعني الأكبر}. فالعدد $3.15$ أرقامه أكثر من
$3.4$ لكنه أصغر منه.
\end{method}

\begin{omfigure}
\begin{tikzpicture}[scale=1.1]
  \draw[-Stealth] (-0.3,0) -- (10.6,0);
  \foreach \i in {0,1,...,10} \draw (\i,-0.08) -- (\i,0.08);
  \node[below, font=\small] at (0,-0.1) {$3$};
  \node[below, font=\small] at (5,-0.1) {$3.5$};
  \node[below, font=\small] at (10,-0.1) {$4$};
  \fill[omThm] (1.5,0) circle (2pt) node[above, font=\small] {$3.15$};
  \fill[omDef] (4,0) circle (2pt) node[above, font=\small] {$3.4$};
  \fill[omMeth] (9,0) circle (2pt) node[above, font=\small] {$3.9$};
\end{tikzpicture}

{\small تكبير مستقيم الأعداد بين $3$ و $4$: كل تدريجة عُشر
واحد. والنقطة $3.15$ تقع في منتصف ما بين $3.1$ و $3.2$.}
\end{omfigure}

\section{الجمع والطرح والضرب}

\begin{method}[الحساب بالأعمدة مع الأعداد العشرية]\label{met:g6:decimals:columns}
لجمع \omterm{def:g6:decimals:places}{الأعداد العشرية} أو طرحها، حاذِ \emph{النقاط العشرية}
(فتقع الآحاد تحت الآحاد، والأعشار تحت الأعشار)، مع الحشو بأصفار
نهائية عند اللزوم؛ ثم احسب كما في الأعداد الصحيحة، وضع نقطة
النتيجة تحت النقطتين الأخريين.
\end{method}

\begin{example}\label{ex:g6:decimals:addsub}
احسب $13.7 + 2.85$، بمحاذاة النقاط ($13.7 = 13.70$):
\[
\begin{array}{r}
1\,3.7\,0 \\
+\ \ \,2.8\,5 \\
\hline
1\,6.5\,5
\end{array}
\]
واحسب $6 - 2.35$ (اُكتب $6 = 6.00$):
$6.00 - 2.35 = 3.65$. والتحقّق: $3.65 + 2.35 = 6$.
\end{example}

\begin{example}[ضرب الأعداد العشرية]\label{ex:g6:decimals:mult}
لحساب $2.3 \times 1.4$: اِضرب كما في الأعداد الصحيحة،
$23 \times 14 = 322$؛ ثم عُدّ الأرقام العشرية في العاملين
($1 + 1 = 2$)، وضع النقطة بحيث يكون في النتيجة هذا العدد منها:
$2.3 \times 1.4 = 3.22$. وتحقّق من رتبة المقدار: \omterm{def:g2:mult:def}{فالجداء} $2.3 \times 1.4$
يجب أن يزيد قليلًا على $2.3 \times 1 = 2.3$ --- والعدد $3.22$ كذلك.
\end{example}

\section{الضرب والقسمة على 10 و100 و1000}

\begin{proposition}[انتقال النقطة]\label{prop:g6:decimals:shift}
ضرب \omterm{def:g6:decimals:places}{عدد عشري} في $10$ أو $100$ أو $1000$ ينقل نقطته
العشرية $1$ أو $2$ أو $3$ مراتب إلى \emph{اليمين}؛ والقسمة تنقلها إلى
\emph{اليسار} (مع الحشو بالأصفار عند اللزوم):
\[
3.75 \times 100 = 375, \qquad
42.1 \div 1000 = 0.0421 .
\]
\end{proposition}

\begin{proof}
الضرب في $10$ يجعل كل رقم يساوي عشرة أمثال ما كان: فتصير الأعشار
آحادًا، والآحاد عشرات، وهكذا --- فينتقل كل رقم عمودًا
واحدًا إلى اليسار في جدول المراتب، وذلك هو نفسه انتقال
النقطة مرتبة واحدة إلى اليمين. والقسمة تعكس ذلك.
\end{proof}

\begin{example}[وحدات القياس]\label{ex:g6:decimals:units}
تحويل الوحدات هو هذه اللعبة بعينها: فبما أن $1$ م $= 100$ سم،
\[
3.75 \text{ م} = 3.75 \times 100 \text{ سم} = 375 \text{ سم},
\qquad
42 \text{ مم} = 42 \div 10 \text{ سم} = 4.2 \text{ سم}.
\]
\end{example}

\section{التقريب}

\begin{definition}[التقريب]\label{def:g6:decimals:rounding}
\emph{تقريب}\index{تقريب} عدد إلى الوحدة (أو العُشر أو
الجزء من مئة، \dots) هو أقرب عدد بتلك الدقة. اُنظر إلى
الرقم التالي: فإن كان $0,1,2,3,4$ قرِّب إلى الأسفل (أبقِ)؛ وإن كان
$5,6,7,8,9$ قرِّب إلى الأعلى.
\end{definition}

\begin{example}
\omterm{def:g6:decimals:rounding}{تقريب} $7.38$ إلى الوحدة هو $7$ (فالرقم التالي $3$: أبقِ)؛ وتقريبه إلى
العُشر هو $7.4$ (فالرقم التالي $8$: قرِّب إلى الأعلى). والثمن
$4.996$ مقرَّبًا إلى الجزء من مئة هو $5.00$ --- \omterm{def:g6:decimals:rounding}{فالتقريب} قد يغيّر
كل الأرقام!
\end{example}

\section{تمارين}

\begin{exercise}[$\star$]\label{exo:g6:decimals:1}
اُكتب على صورة \omterm{def:g6:decimals:places}{عدد عشري}: $5 + \dfrac{3}{10} + \dfrac{7}{100}$؛
\quad $\dfrac{9}{10}$؛ \quad $12 + \dfrac{4}{1000}$؛
\quad «ثماني وحدات وخمسة أجزاء من مئة».
\end{exercise}

\begin{exercise}[$\star$]\label{exo:g6:decimals:2}
في $86.354$: ما رقم الأعشار؟ وما رقم الأجزاء من مئة؟ وماذا
يعدّ الرقم $8$؟ واُكتب هذا العدد كما في
\cref{def:g6:decimals:places}.
\end{exercise}

\begin{exercise}[$\star$]\label{exo:g6:decimals:3}
اُنقل وأكمل بالإشارة $<$ أو $>$ أو $=$:
\[
5.3 \;?\; 5.29, \qquad
0.7 \;?\; 0.70, \qquad
2.09 \;?\; 2.9, \qquad
14.5 \;?\; 14.49 .
\]
\end{exercise}

\begin{exercise}[$\star$]\label{exo:g6:decimals:4}
رتّب من الأصغر إلى الأكبر:
$4.2$؛ \ $4.05$؛ \ $4.51$؛ \ $4.15$؛ \ $4.5$.
\end{exercise}

\begin{exercise}[$\star$]\label{exo:g6:decimals:5}
ما \omterm{def:g6:decimals:places}{الأعداد العشرية} الموافقة للنقاط $A$ و $B$ و $C$ على
مستقيم مدرَّج بالأعشار، إذا كانت $A$ على بُعد $3$ تدريجات بعد $6$،
و $B$ على بُعد $7$ تدريجات بعد $6$، و $C$ على بُعد $2$ من التدريجات بعد $7$؟
\end{exercise}

\begin{exercise}[$\star$]\label{exo:g6:decimals:6}
احسب بالأعمدة: $45.8 + 7.65$؛ \quad $23.4 - 8.72$؛ \quad
$5.6 \times 2.4$ (وعُدّ الأرقام العشرية!).
\end{exercise}

\begin{exercise}[$\star$]\label{exo:g6:decimals:7}
احسب دون أي عمل مكتوب:
\[
6.42 \times 10, \qquad
0.35 \times 1000, \qquad
78.1 \div 100, \qquad
5 \div 1000 .
\]
\end{exercise}

\begin{exercise}[$\star$]\label{exo:g6:decimals:8}
حوِّل: $2.4$ م إلى سم؛ \quad $370$ غ إلى كغ؛ \quad $0.85$ كم إلى
م؛ \quad $56$ مم إلى م.
\end{exercise}

\begin{exercise}[$\star$]\label{exo:g6:decimals:9}
قرِّب $23.867$: إلى الوحدة؛ وإلى العُشر؛ وإلى الجزء من مئة. وقرِّب
$9.97$ إلى العُشر.
\end{exercise}

\begin{exercise}[$\star\star$]\label{exo:g6:decimals:10}
ثمن الرغيف $1.15$. تشتري هناء ثلاثة أرغفة وتدفع بورقة
$5$. اُكتب الحسابين اللازمين، وأعطِ \omterm{def:g3:division:remainder}{الباقي} لها.
\end{exercise}

\begin{exercise}[$\star\star$]\label{exo:g6:decimals:11}
جِد عددًا عشريًّا يقع تمامًا بين $7.4$ و $7.5$؛ ثم بين
$3.99$ و $4$. وكم عدد هذه الأعداد في كل حالة؟
\end{exercise}

\begin{exercise}[$\star\star\star$]\label{exo:g6:decimals:12}
مستعملًا كلًّا من الأرقام $2$ و $5$ و $8$ مرة واحدة بالضبط \omterm{def:g4:decimals:point}{ونقطة
عشرية} واحدة، اُكتب أكبر عدد ممكن، ثم أصغر عدد
ممكن. (والأعداد مثل $.58$ غير مسموحة: فالجزء الصحيح يجب أن يحتوي
رقمًا واحدًا على الأقل.)
\end{exercise}

\section{مسألة: العدد الذي يلي 3 مباشرة لا وجود له}

\begin{problem}[{مسألة نهاية الأسبوع --- بين أي عددين عشريين
يوجد دائمًا عدد آخر: التكبير على مستقيم الأعداد}]
\label{pb:g6:decimals:1}
على سلّم الأعداد الصحيحة، لكل عدد جار ملاصق:
فبعد $7$ مباشرة يأتي $8$، ولا يسكن بينهما شيء.
أما \omterm{def:g6:decimals:places}{الأعداد العشرية} فعالم مختلف تمامًا. فما
العدد الذي يلي $3$ مباشرة؟ أهو $3.1$؟ أم $3.01$؟ أم $3.001$؟ تطوّر هذه
المسألة تقنية \emph{التكبير} على مستقيم الأعداد
(مواصلةً لما في \cref{exo:g6:decimals:11}) وتبلغ نتيجة
مشهورة: العدد الذي يلي $3$ مباشرة \emph{لا وجود له} ---
فبين أي عددين عشريين، مهما تقاربا، يبقى دائمًا
متسع لغيرهما.

\medskip
\noindent\textbf{الجزء الأول --- التكبير.}
\begin{enumerate}
  \item أيّهما أكبر، $2.999$ أم $3$؟ احسب \omterm{ex:g1:subtraction:difference}{الفرق}
        بينهما.
  \item اُرسم مستقيم أعداد من $7.4$ إلى $7.5$، مدرَّجًا
        بالأجزاء من مئة، وضع عليه $7.42$ و $7.45$ و $7.48$.
        واُذكر \emph{كل} الأعداد ذات الرقمين العشريين التي
        تقع تمامًا بين $7.4$ و $7.5$. فكم عددها؟
  \item وكبِّر من جديد: بين $7.44$ و $7.45$، اُذكر كل
        الأعداد ذات الأرقام العشرية الثلاثة. فكم عددها
        هذه المرة؟
  \item وبالفكرة نفسها، اِشرح كيف تعدّ الأعداد
        ذات الأرقام العشرية الثلاثة بالضبط الواقعة تمامًا بين
        $7.4$ و $7.5$، وأعطِ ذلك العدد.
  \item صِف الوصفة الكامنة وراء الأسئلة 2--4 (أي
        \emph{التكبير}): إذا أُعطي عددان يبدوان
        متجاورين، مثل $5.67$ و $5.68$، فكيف تكشف كتابة
        رقم عشري واحد إضافي دائمًا عن عدد يقع تمامًا
        بينهما؟ وطبّق وصفتك على $5.67$ و $5.68$،
        ثم على $0.1999$ و $0.2$.
\end{enumerate}

\medskip
\noindent\textbf{الجزء الثاني --- الجار المفقود.}
\begin{enumerate}[resume]
  \item يدّعي عمر: «العدد الذي يلي $3$ مباشرة هو $3.1$». فأثبت
        خطأه بتسمية عدد يقع تمامًا بين $3$ و
        $3.1$. ثم يتراجع عمر إلى $3.01$، ثم إلى $3.001$. فاغلب كل
        مرشّح من مرشّحيه.
  \item واِشرح لماذا لا يستطيع \emph{أحد} أن يغلبك في هذه اللعبة:
        فمهما اقترح عمر من عدد أكبر تمامًا من $3$،
        تنتج وصفة التكبير في السؤال 5 عددًا يقع تمامًا
        بين $3$ واقتراحه. فماذا يثبت ذلك عن
        «العدد الذي يلي $3$ مباشرة»؟
  \item والآن الجهة الأخرى: يبحث عمر عن العدد الذي
        \emph{يسبق} $3$ مباشرة فيجرّب $2.9$، ثم $2.99$، ثم
        $2.999$. فاغلب مرشّحيه الثلاثة. ثم احسب
        $3 - 2.999$، وصِف (دون حساب بالأعمدة)
        \omterm{ex:g1:subtraction:difference}{الفرق} بين $3$ والعدد المكتوب برقم
        $2$ ونقطة وعشرين رقمًا من $9$.
  \item ولماذا لا ينجح شيء من هذا في الأعداد \emph{الصحيحة}؟
        اِشرح في جملة أو جملتين لماذا لا يوجد عدد
        صحيح يقع تمامًا بين $7$ و $8$، مع أن هناك
        أعدادًا عشرية كثيرة هناك.
  \item وأُعلن طول بأنه «$3.7$ سم، مقرَّبًا إلى
        العُشر» (\cref{def:g6:decimals:rounding}). أعطِ
        أصغر طول يقرَّب إلى $3.7$ سم، واِشرح لماذا
        لا يقرَّب طول قدره $3.75$ سم إلى $3.7$ ---
        فالطول الحقيقي إذن $3.65$ سم على الأقل وأصغر تمامًا
        من $3.75$ سم.
\end{enumerate}

\medskip
\noindent\textbf{الجزء الثالث --- ألعاب بالأرقام والنقط.}
\begin{enumerate}[resume]
  \item رتّب من الأصغر إلى الأكبر: $0.6$؛ \ $0.58$؛ \
        $0.123$؛ \ $0.0999$. ثم اِشرح لزميل، في جملة
        واحدة، لماذا $0.0999 < 0.6$ مع أنه مكتوب
        بأرقام أكثر (\cref{met:g6:decimals:compare}).
  \item والأثمان في متجر لها دائمًا رقمان عشريان بالضبط.
        فهل يوجد \emph{ثمن} يقع تمامًا بين $4.99$ و $5$؟
        قارن بالسؤال 9: فما المشترك بين الأثمان والأعداد
        الصحيحة؟
  \item ومستعملًا كلًّا من الأرقام $2$ و $5$ و $8$ مرة واحدة بالضبط
        \omterm{def:g4:decimals:point}{ونقطة عشرية} واحدة (والجزء الصحيح ليس فارغًا، كما في
        \cref{exo:g6:decimals:12})، جِد العدد الأقرب إلى
        $6$. وبرّر بحساب مسافة أفضل مرشّحيك
        إلى $6$.
  \item وأشهر أعداد الرياضيات، $\pi$، يحقق
        $3.141 < \pi < 3.142$. فسمِّ عددًا عشريًّا يقع تمامًا
        بين $3.141$ و $3.142$؛ ثم عددًا يقع تمامًا بين
        $3.1415$ و $3.1416$. (والرياضيون يحصرون $\pi$
        بهذه الطريقة بالضبط --- فكل رقم عشري جديد تكبير
        إضافي. وستلقى $\pi$ في العمل في
        \cref{ch:g6:measure}.)
  \item والخاتمة: اِشرح لماذا يوجد \emph{من \omterm{def:g6:decimals:places}{الأعداد
        العشرية} بين $0$ و $1$ أكثر من أي عدد تستطيع أن
        تسمّيه}. (فكم عددًا ينتج التكبير الواحد؟ وهل يمكن
        للتكبير أن يتوقف؟)
\end{enumerate}
\end{problem}
